% Part: computability
% Chapter: computability-theory
% Section: no-universal-function

\documentclass[../../../include/open-logic-section]{subfiles}

\begin{document}

\olfileid{cmp}{thy}{nou}
\olsection{تابع محاسبه‌پذیر جهانی وجود ندارد}

هرچند تابعی محاسبه‌پذیر جزئی وجود دارد که برای همهٔ توابع محاسبه‌پذیر
جزئی جهانی است، هیچ تابع محاسبه‌پذیر تامی وجود ندارد که برای همهٔ
توابع محاسبه‌پذیر تام جهانی باشد.

\begin{thm}
\ollabel{thm:no-univ}
هیچ تابع محاسبه‌پذیر جهانی‌ای وجود ندارد. به بیان دیگر، هر تابع
$\fn{Un}'(k,x)$ با این ویژگی که برای هر تابع محاسبه‌پذیر تام~$f(x)$،
عدد طبیعی~$k$ای وجود دارد چنان‌که برای هر~$x$،
$f(x)=\fn{Un}'(k,x)$، محاسبه‌پذیر نیست.
\end{thm}

\begin{proof}
اثبات قطری‌سازی ساده‌ای است: اگر~$\fn{Un}'(k,x)$ تام و محاسبه‌پذیر
بود، آنگاه
\[
d(x) = \fn{Un}'(x, x) + 1
\]
نیز تام و محاسبه‌پذیر می‌بود. اما بنا بر تعریف، $d(k)$ با
$\fn{Un}'(k,k)$ برابر نیست. پس برای هر~$k$، مقادیر~$d(x)$ و
$\fn{Un}'(k,x)$ دست‌کم برای یک~$x$، یعنی $x=k$، متفاوت‌اند.
\end{proof}

\begin{explain}
\olref[uni]{thm:univ-comp} در بالا نشان می‌دهد می‌توانیم از این استدلال
قطری‌سازی بگریزیم، اما تنها به بهای اجازه‌دادن به جزئی‌بودن تابع جهانی.
یعنی~$\fn{Un}$ برای توابع محاسبه‌پذیر تام جهانی است، اما خودش تام
نیست. استدلال قطری‌سازی در حالت جزئی کار نمی‌کند.
\end{explain}

\begin{prob}
  برای فهم اینکه چرا استدلال قطری‌سازی در اثبات
  \olref{thm:no-univ} در حالت جزئی کار نمی‌کند، تابع
  $f(x) \simeq \fn{Un}(x,x)+1$ را در نظر بگیرید. آیا محاسبه‌پذیر جزئی
  است؟ اگر چنین باشد، نمایه‌ای چون~$e$ دارد؛ یعنی
  $f(x) \simeq \fn{Un}(e,x)$. دربارهٔ~$f(e)$ چه می‌توانید گفت؟
\end{prob}


\end{document}
